% ==========================================================================
% SCRIPT D'ORAL — Battle Météorage
% Document à lire en parallèle de la présentation PowerPoint.
% Compilable avec pdflatex sur Overleaf.
% ==========================================================================
\documentclass[11pt,a4paper]{article}

% ------------------ PACKAGES -----------------------------------
\usepackage[utf8]{inputenc}
\usepackage[T1]{fontenc}
\usepackage[french]{babel}
\usepackage{lmodern}
\usepackage[a4paper,margin=2.3cm]{geometry}

\usepackage{amsmath,amssymb,amsthm,bm}
\usepackage{graphicx}
\usepackage{xcolor}
\usepackage{hyperref}
\usepackage{booktabs}
\usepackage{enumitem}
\usepackage{listings}
\usepackage{caption}
\usepackage{float}
\usepackage{tcolorbox}
\usepackage{multirow}
\usepackage{array}
\usepackage{fancyhdr}

% Pas d'indentation de paragraphe (evite le decalage des tcolorbox)
\setlength{\parindent}{0pt}
\setlength{\parskip}{6pt plus 1pt}

\hypersetup{
    colorlinks=true,
    linkcolor=blue!60!black,
    urlcolor=blue!60!black,
    citecolor=blue!60!black
}

% ------------------ STYLE CODE PYTHON --------------------------
\definecolor{codebg}{rgb}{0.97,0.97,0.97}
\definecolor{codekw}{rgb}{0.13,0.29,0.53}
\definecolor{codestr}{rgb}{0.50,0.20,0.00}
\definecolor{codecom}{rgb}{0.25,0.50,0.25}

\lstdefinestyle{py}{
    language=Python,
    backgroundcolor=\color{codebg},
    basicstyle=\ttfamily\footnotesize,
    keywordstyle=\color{codekw}\bfseries,
    stringstyle=\color{codestr},
    commentstyle=\color{codecom}\itshape,
    showstringspaces=false,
    breaklines=true,
    frame=single,
    rulecolor=\color{gray!40},
    numbers=left,
    numberstyle=\tiny\color{gray},
    numbersep=6pt,
    captionpos=b,
    tabsize=4,
    columns=fullflexible
}
\lstset{style=py}

% ------------------ HEADERS ------------------------------------
\pagestyle{fancy}
\fancyhf{}
\fancyhead[L]{Battle Météorage}
\fancyhead[R]{Script d'oral}
\fancyfoot[C]{\thepage}

% ------------------ TITRE --------------------------------------
\title{
    \textbf{Battle Météorage} \\[0.4em]
    \large Script complet pour l'oral \\
    \small Weibull AFT calibré + évaluation CG / 3 km sur eval 2023-25
}
\author{}
\date{\today}

% ==========================================================================
\begin{document}
\maketitle

\tableofcontents
\newpage

% =========================================================================
\section{Introduction (à dire en début d'oral, $\approx$ 1 min)}

\noindent\begin{tcolorbox}[colback=blue!5,colframe=blue!50!black,title=\textbf{Phrase d'accroche}]
\og Bonjour. Nous avons travaillé sur le défi Battle Météorage : peut-on raccourcir la règle des 30 minutes
d'alerte orage actuellement appliquée par Météorage, sans augmenter le risque pour le personnel au sol des aéroports ? Notre réponse : un modèle Weibull AFT calibré, validé par une vérification empirique indépendante,
qui adapte le temps d'attente au contexte de chaque éclair. \fg{}
\end{tcolorbox}

\paragraph{Le contexte.}
Quand un éclair \textbf{nuage-sol} (CG, pour \emph{cloud-to-ground}) tombe dans un rayon de \textbf{20 km} d'un
aéroport, Météorage déclenche une alerte. L'aéroport interrompt ses opérations au sol (atterrissages, ravitaillements,
manutention). L'alerte est levée \textbf{30 minutes après le dernier CG observé}. Chaque nouveau CG remet le
compteur à zéro.

\paragraph{Les enjeux.} Cette règle est aveugle au contexte : tous les orages reçoivent le même délai forfaitaire.
\begin{itemize}[leftmargin=*]
    \item \textbf{Trop court} = risque de foudroiement du personnel ou d'un aéronef au sol $\rightarrow$ inacceptable.
    \item \textbf{Trop long} = retards de vols et pertes opérationnelles $\rightarrow$ coûteux.
\end{itemize}

\paragraph{Notre question.} Sachant qu'il n'y a pas eu de CG depuis $t$ minutes, quelle est la probabilité qu'un
nouveau CG \textbf{dangereux} tombe encore ? Si cette probabilité descend sous un seuil acceptable
$\alpha$ (par exemple 5 \%) au bout de $T < 30$ minutes, on peut remplacer 30 par $T$ et faire gagner du temps
aux opérations.

\paragraph{Notre approche en une phrase.}
\textbf{Un modèle Weibull AFT entraîné par maximum de vraisemblance sur les inter-arrivées CG, calibré
empiriquement, et validé indépendamment par une analyse non-paramétrique.}

% =========================================================================
\section{Les données (à dire $\approx$ 1 min)}

\paragraph{Fichier source.} \texttt{segment\_alerts\_all\_airports\_train.csv}, 507 071 lignes,
période 2016-2022. Une ligne = un éclair détecté dans 30 km d'un aéroport.

\paragraph{Volumétrie après filtrage.} On garde uniquement les CG (\texttt{icloud=False}), dans la zone d'alerte
(\texttt{airport\_alert\_id} non nul), et on exclut Pise 2016 (système d'enregistrement différent) :

\begin{table}[H]
\centering
\begin{tabular}{lrr}
\toprule
\textbf{Aéroport} & \textbf{CG} & \textbf{Alertes} \\
\midrule
Ajaccio  & 10\,647 & 530 \\
Bastia   & 13\,742 & 532 \\
Biarritz &  9\,911 & 590 \\
Nantes   &  4\,378 & 206 \\
Pise     & 15\,792 & 542 \\
\midrule
\textbf{TOTAL} & \textbf{54\,470} & \textbf{2\,400} \\
\bottomrule
\end{tabular}
\end{table}

\paragraph{Structure clé : la notion d'alerte.}
\noindent\begin{tcolorbox}[colback=orange!5,colframe=orange!60!black,title=\textbf{Définition d'une alerte Météorage}]
\begin{itemize}[leftmargin=*,topsep=2pt]
    \item Un CG dans 20 km \textbf{démarre} une alerte.
    \item Toute l'alerte = série de CG espacés de moins de 30 min.
    \item L'alerte se \textbf{ferme 30 min après le dernier CG}.
    \item Un CG plus tard $=$ \textbf{nouvelle alerte} (nouvel \texttt{airport\_alert\_id}).
\end{itemize}
\end{tcolorbox}

\paragraph{Implication essentielle (\og censure à droite\fg{}).}
Dans une même alerte, le gap entre deux CG consécutifs est nécessairement inférieur à 30 min (sinon l'alerte
serait fermée). On ne peut donc \textbf{pas observer} de gaps supérieurs à 30 min intra-alerte. C'est ce qu'on
appelle, en analyse de survie, de la \textbf{censure à droite} : on connaît une borne inférieure (30 min) mais
pas la vraie valeur.

\textbf{À retenir pour l'oral :} \og Notre cible est par construction censurée à 30 min. C'est exactement
ce que les modèles de survie comme Weibull AFT savent gérer. \fg{}

% =========================================================================
\section{Le modèle Weibull AFT — théorie et formules (à dire $\approx$ 3 min)}

\subsection{Pourquoi un modèle de survie}

L'analyse de survie est née en médecine (temps jusqu'à un événement : décès, rechute) et s'applique parfaitement
à notre cas : le \og temps jusqu'au prochain éclair \fg{}. Trois propriétés cruciales :

\begin{enumerate}[leftmargin=*]
    \item \textbf{Gestion native de la censure} : pour le dernier CG d'une alerte, on n'observe pas la durée
    jusqu'au prochain CG (l'alerte se ferme). On sait juste qu'aucun CG n'est arrivé pendant au moins 30 min.
    Cette information est intégrée dans la vraisemblance.
    \item \textbf{Modèle paramétrique avec covariables} : on relie le temps caractéristique $\lambda$ aux variables
    explicatives $X$ via $\log \lambda = \beta_0 + \beta^\top X$ (modèle AFT, \textit{Accelerated Failure Time}).
    \item \textbf{Quantiles analytiques} : la formule $T_q = \lambda \cdot (-\ln q)^{1/k}$ donne directement le
    temps d'attente pour un niveau de risque cible.
\end{enumerate}

\subsection{La distribution de Weibull}

\noindent\begin{tcolorbox}[colback=blue!5,colframe=blue!50!black,title=\textbf{Distribution de Weibull}]
\textbf{Fonction de survie} (probabilité d'attendre plus de $t$ minutes)~:
\[
S(t \mid \lambda, k) = \exp\!\left(-(t/\lambda)^k\right)
\]

\textbf{Densité} et \textbf{hazard} (probabilité instantanée)~:
\[
f(t) = \tfrac{k}{\lambda}\,(t/\lambda)^{k-1}\,S(t), \qquad h(t) = \tfrac{f(t)}{S(t)} = \tfrac{k}{\lambda}\,(t/\lambda)^{k-1}
\]
\end{tcolorbox}

\paragraph{Interprétation des paramètres.}
\begin{itemize}[leftmargin=*]
    \item $\lambda > 0$ : \textbf{échelle} (temps caractéristique, en minutes). Si $\lambda = 2$, on attend en moyenne 2 min.
    \item $k > 0$ : \textbf{forme} (sans unité). Gouverne l'évolution du hazard.
\end{itemize}

\begin{table}[H]
\centering
\begin{tabular}{lll}
\toprule
\textbf{Valeur de $k$} & \textbf{Hazard} & \textbf{Sens physique} \\
\midrule
$k < 1$ & décroissant & orage qui s'éteint (cas typique) \\
$k = 1$ & constant & loi exponentielle, sans mémoire \\
$k > 1$ & croissant & vieillissement, événements plus probables avec le temps \\
\bottomrule
\end{tabular}
\end{table}

\textbf{Sur nos données : $k = 0{,}745$.} Le hazard est décroissant, conforme à l'intuition physique des orages.

\subsection{Modèle AFT avec covariables}

Pour personnaliser la prédiction par éclair, on lie $\lambda$ aux 21 variables $X$ :
\[
\log \lambda(X) = \beta_0 + \beta_1 X_1 + \cdots + \beta_{21} X_{21}
\]

Chaque coefficient $\beta_j > 0$ \textbf{allonge} le temps caractéristique. Un $\beta_j < 0$ le \textbf{raccourcit}.

\subsection{Maximum de vraisemblance avec censure}

\noindent\begin{tcolorbox}[colback=blue!5,colframe=blue!50!black,title=\textbf{Maximum de vraisemblance (MLE)}]
On choisit les paramètres ($\beta$, $k$) qui rendent les données observées les plus probables.

\medskip
\textbf{Cas 1 : éclair non-dernier ($\delta_i = 1$).} On a observé exactement le temps $t_i$ jusqu'au prochain éclair. Le modèle doit donner une forte densité en $t_i$~:
\[
\mathcal{L}_i = f(t_i \mid X_i)
\]

\textbf{Cas 2 : dernier éclair ($\delta_i = 0$, censure).} On sait juste qu'aucun éclair n'est arrivé pendant 30~min. Le modèle doit donner une forte survie à 30~min~:
\[
\mathcal{L}_i = S(30 \mid X_i)
\]

\textbf{Log-vraisemblance totale à maximiser :}
\[
\ell(\beta, k) = \sum_{i=1}^{n} \Big[ \delta_i \log f(t_i \mid X_i) + (1 - \delta_i) \log S(30 \mid X_i) \Big]
\]
\end{tcolorbox}

En Python :
\begin{lstlisting}[caption={Fit du Weibull AFT avec lifelines.}]
aft = WeibullAFTFitter(penalizer=0.05)
aft.fit(train_df, duration_col='duration', event_col='event')
# .fit() trouve par optimisation numérique le bêta et le k
# qui maximisent la log-vraisemblance.
\end{lstlisting}

\subsection{Calcul de $T_q$ — temps recommandé}

Pour un niveau de couverture cible $q$ (ex.\ $0{,}99$ pour 1 \% de risque), on cherche $T_q$ tel que $S(T_q \mid X) = q$ :
\[
T_q(X) = \lambda(X) \cdot (-\ln q)^{1/k}
\]

\textbf{À retenir pour l'oral :} \og Le modèle nous donne, pour chaque éclair et son contexte, le quantile
$T_q$ exact : le temps après lequel la probabilité qu'aucun éclair n'arrive est de $q$. \fg{}

% =========================================================================
\section{Feature engineering — 17 variables contextuelles (à dire $\approx$ 1 min)}

On construit 17 variables qui résument le passé de l'orage au moment de chaque éclair. Aucune ne \og voit \fg{}
le futur (pas de fuite).

\begin{table}[H]
\centering
\small
\begin{tabular}{lll}
\toprule
\textbf{Catégorie} & \textbf{Variables (6+3+4+4 = 17)} & \textbf{Intuition} \\
\midrule
\textbf{Rythme}       & \texttt{cg\_rank, time\_since\_start, prev\_gap,} & rythme passé $\rightarrow$ avenir \\
                      & \texttt{rolling\_gap\_3, rolling\_gap\_5, trend\_gap} & \\
\textbf{Distance}     & \texttt{dist\_current, dist\_cum\_mean, dist\_diff} & orage qui s'éloigne ? \\
\textbf{Intensité}    & \texttt{amp\_abs, amp\_cum\_mean,} & polarité et puissance \\
                      & \texttt{is\_negative, pct\_neg\_cum} & \\
\textbf{Saisonnalité} & \texttt{month\_sin, month\_cos,} & encodage cyclique \\
                      & \texttt{hour\_sin, hour\_cos} & \\
\bottomrule
\end{tabular}
\end{table}

\paragraph{Focus sur \texttt{rolling\_gap\_5}.} Moyenne mobile des gaps entre les 5 derniers CG consécutifs.
\begin{itemize}[leftmargin=*]
    \item Petit ($< 1$ min) $\Rightarrow$ orage dense, encore actif $\Rightarrow$ $T_q$ court
    \item Grand ($> 4$ min) $\Rightarrow$ orage qui se calme $\Rightarrow$ $T_q$ long
\end{itemize}

\paragraph{Encodage cyclique pour les variables temporelles :}
$$\texttt{month\_sin} = \sin(2\pi \cdot \texttt{mois} / 12), \quad \texttt{month\_cos} = \cos(2\pi \cdot \texttt{mois} / 12)$$

Avant le fit, toutes les variables continues sont standardisées (moyenne 0, écart-type 1).

% =========================================================================
\section{Pourquoi pas une simple moyenne ? (à dire $\approx$ 1 min)}

C'est la question naturelle : pourquoi ne pas juste prendre la médiane des gaps par aéroport ?

\noindent\begin{tcolorbox}[colback=purple!5,colframe=purple!50!black,title=\textbf{C-index (Concordance index)}]
Mesure la capacité d'un modèle à \textbf{ordonner} deux observations selon leur vraie durée.
On tire 2 obs au hasard ; si le modèle prédit que la 1re durera plus longtemps et c'est effectivement
le cas, on compte 1. Sinon 0. \textbf{C-index $=$ moyenne sur toutes les paires.}
\begin{itemize}[leftmargin=*]
    \item $0{,}50$ $=$ aléatoire (pile/face)
    \item $1{,}00$ $=$ ordre parfait
    \item $\geq 0{,}70$ $=$ bon modèle
\end{itemize}
\end{tcolorbox}

\begin{table}[H]
\centering
\begin{tabular}{lc}
\toprule
\textbf{Approche} & \textbf{C-index test} \\
\midrule
1. Constant (médiane globale)                   & 0,5000 (aléatoire) \\
2. Médiane par aéroport                         & 0,5459 \\
3. P95 empirique par aéroport                   & 0,5424 \\
4. Weibull AFT, dummies aéroport seuls          & 0,5456 \\
5. \textbf{Weibull AFT + 17 variables}          & \textbf{0,7432} \\
\bottomrule
\end{tabular}
\end{table}

\textbf{Les baselines simples plafonnent à $\sim 0{,}55$} : elles distinguent les aéroports mais rien à
l'intérieur. Notre Weibull AFT atteint $0{,}7432$, soit \textbf{+20 points de discrimination}, exclusivement
grâce à l'adaptation au contexte par éclair.

\paragraph{Exemple concret.} À Ajaccio, deux derniers CG du test :
\begin{itemize}[leftmargin=*]
    \item \textbf{Orage A} (encore actif) : \texttt{rolling\_gap\_5} = 2,07 min, dist = 18,9 km. Weibull AFT prédit
    $T_{99} = 18{,}6$ min.
    \item \textbf{Orage B} (qui se calme) : \texttt{rolling\_gap\_5} = 5,53 min, dist = 10,0 km. Weibull AFT prédit
    $T_{99} = 99{,}7$ min.
\end{itemize}
La médiane par aéroport donnerait $0{,}48$ min pour les deux : incapable de distinguer.

% =========================================================================
\section{Calibration empirique robuste (à dire $\approx$ 2 min)}

\subsection{Le piège : sous-calibration}

\noindent\begin{tcolorbox}[colback=red!5,colframe=red!50!black,title=\textbf{Calibration des quantiles}]
Un modèle est \textbf{bien calibré} si la fréquence empirique d'événements correspond à la probabilité annoncée.
Exemple : si on annonce $T_{99}$ (1 \% de risque), seuls 1 \% des observations doivent dépasser ce $T_{99}$.
\end{tcolorbox}

Sur notre test, le Weibull est sous-calibré :

\begin{table}[H]
\centering
\begin{tabular}{lll}
\toprule
\textbf{On annonce} & \textbf{On a vraiment} & \textbf{Écart} \\
\midrule
10 \% risque & 21 \% risque & $\times 2$ \\
5 \% risque  & 14 \% risque & $\times 3$ \\
1 \% risque  & 6 \% risque  & $\times 6$ \\
\bottomrule
\end{tabular}
\end{table}

\textbf{Pourquoi ?} La queue de Weibull décroît comme $\exp(-(t/\lambda)^k)$ — c'est rapide. Mais les vrais longs
silences sont \emph{plus longs} que ce que Weibull prédit. Le modèle est trop optimiste sur les extrêmes.

\subsection{Solution en deux étapes : winsorization $+$ monotonicité}

\paragraph{Étape A — winsorization.}
\noindent\begin{tcolorbox}[colback=green!5,colframe=green!50!black]
\textbf{Winsorization} : on remplace les valeurs au-dessus du percentile 95 par la valeur du P95 lui-même.
Retire l'effet des outliers extrêmes \textbf{sans les supprimer}.
\end{tcolorbox}

\begin{lstlisting}[caption={Winsorization des résidus avant calcul de $c_q$.}]
for q in [0.90, 0.95, 0.99]:
    Tq_pred = aft.predict_percentile(test_evt, p=1.0 - q).values
    residuals = y_real / np.where(Tq_pred > 1e-9, Tq_pred, 1e-9)
    p95 = np.quantile(residuals, 0.95)
    residuals_wins = np.minimum(residuals, p95)
    c_q_winsorized[q] = np.quantile(residuals_wins, q)
\end{lstlisting}

\paragraph{Étape B — monotonicité de $c_q$.}
À $q=99\%$ on obtient $c_{99}^{\text{wins}} = 1{,}13 < c_{95} = 2{,}02$ — ce qui rendrait $T_{99} < T_{95}$,
mathématiquement impossible. On impose :
$$c_q^{\text{final}} = \max(c_q^{\text{winsorisé}},\ c_{q-1}^{\text{final}})$$

\begin{table}[H]
\centering
\begin{tabular}{lccc}
\toprule
$q$ & $c_q^{\text{brut}}$ & $c_q^{\text{winsorisé}}$ & $c_q^{\text{final}}$ \\
\midrule
10 \% & 1,82 & 1,82 & 1,82 \\
5 \%  & 2,02 & 2,02 & 2,02 \\
1 \%  & \textbf{2,82} & \textbf{1,13} & \textbf{2,02} \\
\bottomrule
\end{tabular}
\end{table}

\subsection{Résultat final calibré}

\[
T_q^{\text{calibré}}(X) = c_q \cdot \hat{T}_q^{\text{Weibull}}(X)
\]

\begin{itemize}[leftmargin=*]
    \item À 10 \% de risque : $T = 16$ min (gain de 14 min)
    \item À 5 \% de risque : $T = 26$ min (gain de 4 min)
    \item À 1 \% de risque : $T = 46$ min — au-dessus de 30, on garde la baseline
\end{itemize}

% =========================================================================
\section{Vérification empirique non-paramétrique (à dire $\approx$ 1 min)}

On dispose d'une approche \textbf{complètement non-paramétrique} qui ne fait aucune hypothèse de
distribution. Elle sert de \textbf{vérification indépendante} du Weibull AFT.

\paragraph{Définition.}
Soit une alerte $A$ contenant les CG aux instants $t_1 < t_2 < \cdots < t_n$.
\begin{itemize}[leftmargin=*]
    \item Inter-arrivées : $\Delta_i = t_{i+1} - t_i$, pour $i = 1, \dots, n-1$.
    \item \textbf{Plus grand silence intra-alerte} : $G(A) = \max_{1 \leq i \leq n-1} \Delta_i$
\end{itemize}

Si on avait choisi un temps de maintien $T < G(A)$, on aurait fermé l'alerte \textbf{avant} le CG suivant —
c'est un incident.

\noindent\begin{tcolorbox}[colback=green!5,colframe=green!50!black,title=\textbf{Risque empirique}]
Pour un $T$ candidat~:
\[
\widehat{R}(T) = \tfrac{1}{N} \sum_{k=1}^{N} \mathbb{1}\!\left[G(A_k) > T\right]
\]

On cherche~:
\[
T^*(\alpha) = \min\{\,T : \widehat{R}(T) \leq \alpha\,\}
\]
\end{tcolorbox}

\paragraph{Résultats sur train.}
\begin{table}[H]
\centering
\begin{tabular}{lcc}
\toprule
\textbf{Tolérance $\alpha$} & \textbf{$T^*(\alpha)$} & \textbf{Gain vs 30 min} \\
\midrule
10 \% & 25 min & $+5$ min \\
5 \%  & 28 min & $+2$ min \\
\textbf{1 \%}  & \textbf{30 min} & \textbf{0 min} (= règle Météorage !) \\
\bottomrule
\end{tabular}
\end{table}

\noindent\begin{tcolorbox}[colback=blue!5,colframe=blue!50!black]
\textbf{Découverte} : à $\alpha = 1$ \%, $T^*(1\%) = 30$ min — \emph{exactement} la règle Météorage.
La règle des 30 min n'est pas un choix arbitraire : elle est statistiquement bien calibrée pour 1 \% de risque
résiduel.
\end{tcolorbox}

% =========================================================================
\section{Évaluation sur eval 2023-25 — simulation opérationnelle réelle (à dire $\approx$ 4 min)}

\subsection{Le critère opérationnel}

Le risque opérationnel pour un aéroport se joue à \textbf{moins de 3 km} (zone du personnel
sol, ravitaillement, chargement). Le protocole officiel du Data Battle impose donc le
critère suivant :

\noindent\begin{tcolorbox}[colback=orange!5,colframe=orange!60!black,title=\textbf{Définition de l'incident}]
\textbf{Incident} = un \textbf{CG} (nuage-sol) à $<3$ km de l'aéroport arrivant \emph{après}
la fin d'alerte prédite par notre modèle. Les IC (intra-nuage) ne comptent pas — ils ne
touchent pas le sol.
\end{tcolorbox}

\subsection{La mécanique opérationnelle réelle}

\textbf{À dire} : \og On ne sait jamais en temps réel quand un orage s'arrête. La seule
chose qu'on fait, c'est démarrer un timer après chaque CG. Si le timer expire sans nouveau CG,
on lève l'alerte. Si un nouveau CG arrive avant, on reset. C'est exactement la mécanique
Météorage 30 min, mais avec $T_q$ adapté au contexte. \fg{}

\noindent\begin{tcolorbox}[colback=blue!5,colframe=blue!50!black,title=\textbf{Simulation pas à pas}]
Pour chaque alerte de l'eval, on simule CG par CG :
\begin{enumerate}[topsep=2pt]
    \item À chaque CG observé, le modèle prédit $T_q(X)$ minutes.
    \item Le \textbf{timer effectif} démarre : $\text{timer} = \max(T_q, T_{\min})$ où
    $T_{\min}$ est un \textbf{plancher opérationnel} (l'opérateur ne fait pas confiance
    aveuglément au modèle).
    \item Si un \textbf{nouveau CG arrive avant} l'expiration du timer $\to$ RESET au nouveau CG.
    \item Si le \textbf{timer expire sans nouveau CG} $\to$ LEVER l'alerte à ce moment-là.
\end{enumerate}
Puis on compte les CG <3 km arrivant après notre levée (incidents) et le temps gagné par
rapport à la baseline Météorage (= dernier CG + 30 min).
\end{tcolorbox}

\paragraph{Le plancher $T_{\min}$ — levier opérationnel.}
\begin{itemize}[leftmargin=*]
    \item $T_{\min} = 0$ : on suit pleinement le modèle. Risque maximum, gain maximum.
    \item $T_{\min} = 30$ : on attend toujours $\geq 30$ min après chaque CG = baseline Météorage.
    Risque minimum, gain minimum.
    \item $T_{\min}$ intermédiaire : compromis.
\end{itemize}

\subsection{Résultats globaux — sweep $T_{\min}$ à $q = 0{,}95$}

Simulation sur \textbf{973 alertes} de l'eval 2023-25 (\textbf{385 CG <3 km dangereux}) :

\begin{table}[H]
\centering
\small
\begin{tabular}{ccccc}
\toprule
\textbf{$T_{\min}$} & \textbf{Timer moyen} & \textbf{Gain total} & \textbf{Gain/alerte} & \textbf{Risque CG <3 km} \\
\midrule
 0 min &  4,9 min & 953 h & 58,8 min & 88,05 \% (339/385) — \textit{non déployable} \\
10 min & 12,4 min & 673 h & 41,5 min & 53,25 \% (205/385) \\
15 min & 17,6 min & 505 h & 31,1 min & 33,77 \% (130/385) \\
20 min & 22,0 min & 365 h & 22,5 min & 28,05 \% (108/385) \\
\textbf{25 min} & \textbf{26,2 min} & \textbf{251 h} & \textbf{15,4 min} & \textbf{17,14 \% (66/385)} — \textit{recommandé} \\
30 min & 30,0 min & 136 h &  8,4 min &  8,57 \% (33/385) — \textit{prudent} \\
\midrule
Baseline 30 min Météorage stricte & 30 min & 0 h & 0 min & 0 \% (0/385) — référence \\
\bottomrule
\end{tabular}
\caption{Compromis gain/risque sur eval 2023-25. Le sweet spot opérationnel est à
$T_{\min} = 25$ : 251 h gagnées sur la durée de l'eval, 17 \% de risque CG, timer moyen
de 26 min.}
\end{table}

\textbf{À dire} : \og Le plancher $T_{\min}$ est le levier qu'un opérateur choisit selon sa
tolérance au risque. Sans plancher, le modèle lèverait souvent trop tôt (88 \% de risque).
Avec un plancher de 25 min, on garde 251 heures de gain au prix d'un risque résiduel
de 17 \%. \fg{}

\subsection{Détail par aéroport (à $q = 0{,}95$, $T_{\min} = 25$)}

\begin{table}[H]
\centering
\small
\begin{tabular}{lcccccc}
\toprule
\textbf{Aéroport} & \textbf{Nb alertes} & \textbf{Timer moyen} & \textbf{Gain/alerte} & \textbf{Gain total} & \textbf{CG manqués} & \textbf{Risque CG} \\
\midrule
Ajaccio  & 192 & 26,0 min & 11,5 min & 37 h & 7 / 43  & 16,28 \% \\
Bastia   & 233 & 26,3 min & 15,9 min & 62 h & 9 / 122 &  7,38 \% \\
Biarritz & 217 & 26,4 min & 12,7 min & 46 h & 7 / 55  & 12,73 \% \\
Nantes   &  88 & 25,7 min & 16,0 min & 24 h & 7 / 26  & 26,92 \% \\
Pise     & 243 & 26,1 min & 20,4 min & 83 h & 36/ 139 & 25,90 \% \\
\midrule
\textbf{TOTAL} & \textbf{973} & \textbf{26,2 min} & \textbf{15,4 min} & \textbf{251 h} & \textbf{66 / 385} & \textbf{17,14 \%} \\
\bottomrule
\end{tabular}
\caption{Détail par aéroport. Les profils diffèrent fortement : Pise (orages denses)
gagne le plus (83 h) mais a un risque élevé (26 \%). Bastia gagne 62 h avec seulement
7,4 \% de risque (proportionnellement le meilleur). Nantes a peu d'alertes mais un risque
relatif élevé (7/26 = 27 \%).}
\end{table}

\textbf{À dire} : \og On voit que tous les aéroports ne se comportent pas pareil. Pise
gagne 83 heures avec 26 \% de risque — c'est le plus généreux. Bastia gagne 62 heures
avec seulement 7 \% de risque — proportionnellement le meilleur. Le modèle s'adapte
au profil local de chaque aéroport. \fg{}

\subsection{Compromis $(q, T_{\min})$ — Pareto frontier}

On peut aussi sweep $q$ (calibration du quantile) :

\begin{table}[H]
\centering
\small
\begin{tabular}{ccccc}
\toprule
\textbf{$q$} & \textbf{$T_{\min}$} & \textbf{Gain total} & \textbf{Risque CG} & \textbf{Profil} \\
\midrule
0,85 & 25 & 260 h & 17,40 \% & agressif (T\_q court) \\
0,95 & 25 & 251 h & 17,14 \% & équilibré \textit{(recommandé)} \\
0,99 & 25 & 234 h & 16,10 \% & prudent (T\_q long) \\
\midrule
0,95 & 20 & 365 h & 28,05 \% & plus de gain, plus de risque \\
0,95 & 30 & 136 h &  8,57 \% & plus prudent, moins de gain \\
\bottomrule
\end{tabular}
\caption{Le choix de $q$ a un effet modeste (le plancher $T_{\min}$ domine). Le levier
principal est $T_{\min}$.}
\end{table}

\subsection{Pourquoi le risque n'est pas à 0 \% : limites assumées}

\noindent\begin{tcolorbox}[colback=red!5,colframe=red!50!black,title=\textbf{Honnêteté scientifique}]
La cible 2 \% du protocole officiel n'est pas atteinte. Plusieurs causes :
\begin{enumerate}[topsep=2pt]
    \item \textbf{Notre confidence $S(30|X)$ est imparfaite} : à un CG en milieu d'orage actif,
    le modèle prédit $T_q$ court (orage encore vivant). Si le prochain CG arrive plus tard
    que $T_q$ par chance, on lève prématurément → incident.
    \item \textbf{La définition d'alerte dans l'eval est plus large que la règle stricte 30 min} :
    108 gaps intra-alerte > 30 min existent dans l'eval. Notre simulation lève à $T_{\min} = 30$
    en interprétant strictement, mais Météorage maintient l'alerte ouverte plus longtemps
    dans ces cas. D'où le 8,57 \% résiduel même à $T_{\min} = 30$.
\end{enumerate}

\textbf{Pour atteindre R < 2 \%} : il faudrait entraîner un \textbf{classifieur secondaire dédié}
à la tâche \og ce CG est-il le dernier de l'alerte ? \fg{} avec des features combinant $T_q$,
$S(30|X)$, \texttt{rolling\_gap}, \texttt{cg\_rank}, etc. C'est notre perspective principale
d'amélioration.
\end{tcolorbox}

\subsection{Comparaison avec le protocole officiel Evaluation\_databattle\_meteorage}

Notre simulation reprend exactement la mécanique du protocole officiel :
\begin{itemize}[leftmargin=*,topsep=2pt]
    \item \textbf{Pareil} : 1 décision de lever par alerte, gain = (last\_lightning + 30 min) − levée\_prédite,
    risque = CG <3 km arrivant après levée\_prédite.
    \item \textbf{Différence} : le notebook officiel utilise un fichier \texttt{predictions.csv}
    avec multi-prédictions + sweep $\theta$ sur la confidence. Nous utilisons la mécanique
    temps réel (timer CG par CG) avec un plancher $T_{\min}$. Les deux donnent des résultats
    similaires si la confidence et le plancher sont calibrés équivalents.
\end{itemize}

% =========================================================================
\section{Risque empirique non-paramétrique adapté à 3 km}

Sur l'eval, dans l'esprit de la section 7 :
$$\widehat{R}(T) = \frac{|\{i : \Delta_i > T \text{ ET } \text{dist}(CG_{i+1}) < 3\}|}{N_{\text{CG}<3\text{km}}}$$

\begin{table}[H]
\centering
\begin{tabular}{lcc}
\toprule
\textbf{$\alpha$ visé} & \textbf{$T^*(\alpha)$} & \textbf{Gain vs 30 min} \\
\midrule
10 \% & 5 min  & $+25$ min \\
5 \%  & 9 min  & $+21$ min \\
1 \%  & $>30$ min & non atteignable \\
\bottomrule
\end{tabular}
\end{table}

\textbf{Interprétation à dire} :
\og Cette approche non-paramétrique, indépendante du modèle, confirme qu'un temps fixe inférieur à 30 min
peut être tolérable jusqu'à un certain niveau de risque. Pour atteindre 1 \% de risque sur les CG dangereux,
il faut maintenir la baseline 30 min. \fg{}

% =========================================================================
\section{Démonstration Streamlit (à montrer pendant l'oral $\approx$ 2 min)}

\paragraph{Lancement.}
\begin{lstlisting}[language=bash]
streamlit run app.py
\end{lstlisting}

\paragraph{Parcours à faire devant le jury.}
\begin{enumerate}[leftmargin=*]
    \item Sélectionner $q = 0{,}95$ dans la sidebar.
    \item Cliquer sur \og Lancer l'évaluation \fg{}.
    \item Lire les KPI globaux : 17 037 CG analysés, 37 incidents, gain ${\sim}$ X heures, risque réel 9,61 \%.
    \item Descendre dans \og Les 37 incidents détectés \fg{}, ouvrir le premier : \texttt{Pise\_1319}.
    \item Lire la timeline pas à pas (1 à 5 dans l'app) : \og Le modèle a prédit 2,75 min, mais le prochain
    CG est arrivé 35 min plus tard à 1,87 km. C'est un incident. \fg{}
    \item Montrer la courbe de risque empirique non-paramétrique en bas.
    \item Conclure : \og La démo est entièrement interactive : on peut faire varier $q$ et observer le
    compromis entre gain et risque. \fg{}
\end{enumerate}

% =========================================================================
\section{Limites et perspectives (à dire $\approx$ 30 s)}

\paragraph{Limites assumées.}
\begin{itemize}[leftmargin=*]
    \item Données censurées à 30 min par construction.
    \item Le modèle prédit le temps jusqu'au \emph{prochain} CG, pas le temps jusqu'au \emph{dernier} CG
    de l'orage. C'est ce qui explique la différence entre mesure dynamique ($\sim 10\%$ à $q=0{,}95$) et
    protocole officiel ($\sim 0{,}6$ \%).
    \item Calibration sur le test (en production, prévoir un split train / calibration / test distinct).
    \item Pas de météo externe (vent, pression, humidité).
\end{itemize}

\paragraph{Prochaines étapes.}
\begin{itemize}[leftmargin=*]
    \item Intégrer les données météo via API Météo France.
    \item Re-fit annuel pour absorber la dérive climatique.
    \item Validation A/B en production.
    \item Tester le mode \texttt{ancillary} (paramètre $k$ variable selon $X$).
\end{itemize}

% =========================================================================
\section{Conclusion (à dire $\approx$ 30 s)}

\noindent\begin{tcolorbox}[colback=green!5,colframe=green!60!black,title=\textbf{Synthèse pour la conclusion orale}]
\og Notre Weibull AFT calibré prédit, pour chaque éclair, un temps d'attente $T_q$ adapté
au contexte de l'orage. On l'évalue avec une simulation opérationnelle réelle : timer CG par
CG avec un plancher $T_{\min}$ qui garantit qu'on ne lève jamais avant ce seuil.

Sur eval 2023-25 (973 alertes, 385 CG <3 km), à $q = 0{,}95$ :
\begin{itemize}[topsep=0pt]
    \item \textbf{Point opérationnel recommandé} ($T_{\min} = 25$) : \textbf{251 heures gagnées},
    timer effectif moyen de 26 min, risque CG = \textbf{17 \%} (66 CG manqués / 385).
    \item Point prudent ($T_{\min} = 30$) : 136 h gain, 8,6 \% de risque.
    \item Point agressif ($T_{\min} = 20$) : 365 h gain, 28 \% de risque.
    \item Baseline 30 min Météorage : 0 h / 0 \% (référence parfaite).
\end{itemize}
La démo Streamlit permet de visualiser chaque incident et de faire varier $(q, T_{\min})$
en live. Le code, le rapport et la présentation sont fournis. Merci. \fg{}
\end{tcolorbox}

% =========================================================================
\appendix
\section{Annexe — fichiers livrés}

\begin{itemize}[leftmargin=*]
    \item \texttt{Weibull\_final3km.ipynb} : notebook principal d'évaluation 3 km (mesure dynamique)
    \item \texttt{Evaluation\_databattle\_weibull.ipynb} : reproduction du protocole officiel
    du Data Battle, sweep $\theta$ sur \texttt{predictions.csv}
    \item \texttt{weibull\_final.ipynb} : notebook de base (théorie + calibration sur train)
    \item \texttt{app.py} : démo Streamlit (lancer \texttt{streamlit run app.py})
    \item \texttt{meteorage\_model.py} : module Python réutilisable (load, features, fit)
    \item \texttt{rapport\_final.tex} : rapport long
    \item \texttt{presentation\_oral.pptx} : présentation PowerPoint
    \item \texttt{script\_oral.tex} : ce document
\end{itemize}

\end{document}
